\section{Getting to a language}
This section describes the process, the steps, and the considerations for
designing the risc-v styled typed assembly language suited for static
worst-case execution time analysis.
Only a subset of the idea of the language will be fully fleshed out, and an
even smaller part implemented.
The \say{full} language requires much too complex typing that involves
constraint satisfaction of predicate logic, far out of the scope of this
project, unfortunately.

\subsection{Strays from risc-v}

As with the TALs introduced above, there is little reason to distinguish
between immediate instructions and register instructions, meaning that \subi
and \sub can be treated the same. This is due to the value type that they all
share in some form, defined in some way as
\begin{equation*}
    \begin{array}{llcl}
        \text{value}      & v     & ::= & c \mid r \\
    \end{array}
\end{equation*}
For code generation, these will have to be reintroduced, but this is trivial,
since the value field will either be an immediate value or a register.

\subsection{Stacks and arrays}

This section describes how the array and stack types work with each other.
From the syntax we have

\begin{figure}
    \begin{equation*}
        \begin{array}{llcl}
            \text{region memory types variables} & \mu & & \\
            \text{region memory types}      & M     & ::= & scratch \mid shared \mid \mu \\
            \text{stack types variables}    & \rho  & & \\
            \text{stack types}              & S     & ::= & [] \mid \rho \mid \tau :: S \\
            \text{types}                    & \tau  & ::= & .. \mid
                                                            \tau\ array(x)
                                                            \\
        \end{array}
    \end{equation*}
    \caption{Syntax for lifting types into region memory}
    \label{syn:region-lifted}
\end{figure}

Stack types can be instantiated with

\begin{equation*}
    \infer[(::)]
    { \Phi; \Delta; (R, S) \vdash \subi\ r_d, r_s, imm; I}
    { \Phi; \Delta; (R, S) \vdash r_s : \rho 
    & imm / 4 = n
    & n > 0
    & \Phi; \Delta; (R[r_d : top_0 :: \cdots :: top_n :: \rho], S) \vdash 
    }
\end{equation*}

This of course lacks some expressiveness, since rarely is an extended 

this immediately creates a problem, since only the risc-v stack grows
downwards, and the region stacks (of which allocations happen to). However,
there is a solution in the conceptualization of these stacks.

An example program that stores some values on the stack would look like the
following
\begin{lstlisting}[gobble=4]
    start:  [a0: int(i), a1: int(j), a2: int(k), a3: int(l)]
            subi    sp, sp, 16      ; sp: scratch top :: top :: top :: top :: 'r
            mv      a4, sp
    coerce: ('m)
            [a0: int(i), a1: int(j), a2: int(k), a3: int(l), a4: 'm int array(4)]
            lw      a0, 0(a4)       ; a4: scratch int(i) :: ..
            lw      a1, 4(a4)       ; a4: scratch int(i) :: int(j) :: ..
            lw      a2, 8(a4)       ; a4: scratch int(i) :: int(j) :: int(k) ..
            lw      a3, 12(a4)      ; a4: scratch int(i) :: int(j) :: int(k) :: int(l) :: 'r
\end{lstlisting}
% \textit{dependent types are not making it into the implementation, and all arrays will be statically sized}
\cons is right associative, so $\texttt{scratch top :: top ..} =
\texttt{scratch} (\texttt{top :: top ..})$

\ttop can be coerced into any type, and is used to initialize the array.

A stack $[\tau_0 \cons \cdots \cons \tau_n]$ can be coerced into an array
$\tau\ array(\var m)$ iff $m < n$ and $\tau = tau_i$ for all $0 \leq i < m$

\var{'m} is a type variable that will be unified to type \scratch.

Since the stack pointer points to the stack, and the target is a specific ssd
os system (link) the type of \var{sp} will always be $\scratch \mathrm S$.

Another use is working with the arena allocators. They are simple, and also
work like a stack, but these grow upwards, which is opposite the call stack and
therefore the typing rules would not work.

\begin{lstlisting}[language=C, gobble=4]
    struct memory_region
    {
        uint8_t *start;
        size_t   size;
    };

    struct arena_allocator
    {
        struct memory_region region;
        size_t               allocated;
    };
\end{lstlisting}

\subsubsection{Unitialized values} Since the stack \incr instruction
essentially creates garbage values, we may want to type it, and introduce a
notion of instantiated and uninstantiated types\cite{stal}.

Essentially, the regions type --- while in C is just a pointer --- is a
lifted stack type containing the whole region.
When memory is allocated in an arena, the type of the returned memory is a
coerced slice of the arena.
If an arena has region of type $\ttop_0 \cons \cdots \cons \ttop_n \cons
\emptystack$, then an offset of $k$ words gives the stack $\ttop_k \cons \cdots
\cons \ttop_n \cons \emptystack$, which can be coerced into an array of fitting
type and size.

Since the language --- like other TALs --- only work on word sized types but
operations on memory in risc-v (like moving the stack pointer) is not aligned,
we want to constrain the syntax to only allow word-size operations.
Only stacks can change size (through arithmetic operations on a stack pointer)
so it suffices to introduce instructions \incr and \decr.
They may be best explained through their code generation, shown in figure \ref{fig:incr-codegen}
\begin{figure}
    \begin{equation*}
        \begin{array}{lcl}
            \incr\ r_d,\ r_s,\ v
            & \mapsto
            & \left\{
                \begin{array}{lcl}
                    \begin{array}{llll}
                        \mul, \quad    & r_d,\ & v,\   & wordsize \\
                        \sub,           & r_d,  & r_s,  & rd
                    \end{array}
                    & \text{iff}
                    & \Phi; \Delta; R \vdash v : \text{is register}
                    \\
                    \\
                    \begin{array}{llll}
                        \subi, \quad    & r_d,  & r_s,  & v \cdot wordsize
                    \end{array}
                    & \text{iff}
                    & \Phi; \Delta; R \vdash v : \text{is immediate value}
                \end{array}
            \right.
            \\
            \\
            \decr\ r_d,\ r_s,\ v
            & \mapsto
            & \left\{
                \begin{array}{lcl}
                    \begin{array}{llll}
                        \mul, \quad    & r_d,\ & v,\   & wordsize \\
                        \add,           & r_d,  & r_s,  & rd
                    \end{array}
                    & \text{iff}
                    & \Phi; \Delta; R \vdash v : \text{is register}
                    \\
                    \\
                    \begin{array}{llll}
                        \addi, \quad    & r_d,  & r_s,  & v \cdot wordsize
                    \end{array}
                    & \text{iff}
                    & \Phi; \Delta; R \vdash v : \text{is immediate value}
                \end{array}
            \right.
        \end{array}
    \end{equation*}
    \caption{Code generation for \incr and \decr}
    \label{fig:incr-codegen}
\end{figure}


\subsection{Types}
    As discussed, we want to lift types into memory regions as in figure
    \ref{syn:region-lifted}.
    The evaluation will be as in TAL-0

\subsection{Abstract machine}
    Unlike other TAL abstract machines, I am omitting the stack, since it will
    be used as a normal register (\var{sp}).

\subsection{Dependent types}

Because of the restrictions on the implementation domain, we can restrict the
dependent types to contain fewer index propositions.

% Dependent types are a mess, since

\begin{lstlisting}
    L0:     {n: int | n > 0}
            [a0: int(n)]
            muli    a1, a0, 4       ; a1: n * 4
            sub     sp, sp, a1      ; sp: top_0 :: ... :: top_n :: sp
            mv      a1, a0
            mv      a0, sp
    L1:     ('m)
            {n: int | n > 0}
            [a0: 'm int array(n), a1: int(n)]
\end{lstlisting}
Is a valid program. The two \var{n}s are locally scoped, so they share no
information, however they should represent the same value in this case.
Since \var n is not known, only that it is a natural non-zero number, it is not
known how many words are pushed to the stack, therefore its type depends on
\var n.

However, since the lack of control flow, loops do not exist and the language
turns novel. Then there is no way to store \var n words, and therefore no
reason to be able to allocate a variable amount of words. The language reduces
to constant expressions. In future work, reintroducing control flow is a must,
and as discussed previously (TODO: HOPEFULLY) dependent types can be used to
infer some loop iterations. For this, we would sorely miss type variables in
index expressions.


this reduces stack type rules to
\begin{equation*}
    \infer[(\sub)]
    {\Phi;\ \Delta;\ R \vdash \sub\ r_d,\ r_s,\ v}
    {}
\end{equation*}


\subsection{Coercion}

\subsection{Syntax}
\begin{figure}
    \begin{equation*}
        \begin{array}{llcl}
            \text{type variables}           & \alpha \\
            \text{region memory types variables} & \mu \\
            \text{region memory types}      & M     & ::= & scratch \mid shared \mid \mu \\
            \text{stack types variables}    & \rho  & & \\
            \text{stack types}              & S     & ::= & [] \mid \rho \mid \tau :: S \\
            \text{Collection types}         & C     & ::= & M\ \tau\ array(x)
                                                        \mid M\ S
                                                        \\
            \text{types}                    & \tau  & ::= & \alpha
                                                        \mid top
                                                        \mid int
                                                        \mid C
                                                        \\
            \text{type variable contexts}   & \Delta & ::= & \cdot_\mathrm{tv} 
                                                        \mid \Delta, \mu 
                                                        \mid \Delta, \rho 
                                                        \mid \Delta, \alpha
                                                        \\
        \end{array}
    \end{equation*}
    \caption{Syntax for a typed assembly language}
    \label{syn:my-tal}
\end{figure}

An example program could look like figure \ref

\begin{figure}
    \begin{lstlisting}[gobble=8]
        start:  [a0: int, a1: int, a2: int, a3: int]
                subi    sp, sp, 16      ; sp: scratch top :: top :: top :: top :: 'r
                mv      a4, sp
        coerce: ('m)
                [a0: int, a1: int, a2: int, a3: int, a4: 'm int array(4)]
                lw      a0, 0(a4)       ; a4: scratch int :: ..
                lw      a1, 4(a4)       ; a4: scratch int :: int :: ..
                lw      a2, 8(a4)       ; a4: scratch int :: int :: int ..
                lw      a3, 12(a4)      ; a4: scratch int :: int :: int :: int :: 'r
    \end{lstlisting}
    \caption{Example program}
    \caption{code:mtal}
\end{figure}
